% Chapitre 5 : Échantillonnage et TCL

\section{Population, échantillon, statistique}

\begin{Def}
La \textbf{population} est l'ensemble complet sur lequel on voudrait de l'information.
\end{Def}

\begin{Def}
Un \textbf{échantillon} est un ensemble fini d'observations prélevées dans la population.
\end{Def}

\begin{Ex}
\textbf{Exemples}
\begin{itemize}
    \item Population : toutes les éprouvettes produites par une centrale pendant une campagne.
    \item Échantillon : les $n=30$ éprouvettes contrôlées cette semaine.
\end{itemize}
\end{Ex}

\begin{Def}
Une \textbf{statistique} est une quantité calculée à partir de l'échantillon : moyenne, proportion, variance empirique, maximum, etc.
\end{Def}

\section{Moyenne et proportion empiriques}

\begin{Def}
Si l'on observe $x_1,\dots,x_n$, la \textbf{moyenne empirique} est
$$
\overline{x}=\frac{x_1+\cdots+x_n}{n}.
$$
\end{Def}

\begin{Def}
Si l'on compte un nombre $k$ de succès dans un échantillon de taille $n$, la \textbf{proportion empirique} est
$$
\widehat{p}=\frac{k}{n}.
$$
\end{Def}

\begin{Rmq}
La moyenne empirique et la proportion empirique sont des observations calculées sur l'échantillon. Elles servent à approcher la moyenne ou la proportion de la population.
\end{Rmq}

\section{Dispersion dans l'échantillon}

\begin{Def}
La \textbf{variance empirique} est
$$
s^2=\frac{1}{n}\sum_{i=1}^n (x_i-\overline{x})^2.
$$
\end{Def}

\begin{Def}
L'\textbf{écart-type empirique} est $s=\sqrt{s^2}$.
\end{Def}

\begin{Rmq}
Dans ce cours, l'objectif principal de $s$ est d'évaluer la dispersion d'un échantillon et de comprendre la précision d'une moyenne.
\end{Rmq}

\section{Moyenne d'échantillon comme variable aléatoire}

\begin{Prop}
Si $X_1,\dots,X_n$ sont indépendantes et de même loi, de moyenne $\mu$ et de variance $\sigma^2$, alors la moyenne d'échantillon
$$
\overline{X}=\frac{X_1+\cdots+X_n}{n}
$$
verifie
$$
\mathbb{E}[\overline{X}]=\mu,
\qquad
\Var(\overline{X})=\frac{\sigma^2}{n}.
$$
\end{Prop}

\begin{Rmq}
Quand la taille d'échantillon augmente, la variance de la moyenne diminue. Une moyenne calculée sur beaucoup d'observations est donc plus stable.
\end{Rmq}

\section{Théorème central limite}

\begin{Prop}[Théorème central limite]
Si $X_1,\dots,X_n$ sont indépendantes et de même loi, de moyenne $\mu$ et d'écart-type $\sigma$, alors pour $n$ grand :
$$
\frac{\overline{X}-\mu}{\sigma/\sqrt{n}}
$$
est approximativement distribuée selon la loi normale centrée réduite $\mathcal{N}(0,1)$.
\end{Prop}

\begin{Rmq}
Le théorème central limite explique pourquoi la loi normale apparaît si souvent en statistique, même lorsque la population de départ n'est pas normale.
\end{Rmq}

\section{Cas d'une proportion}

\begin{Prop}
Si l'on observe $n$ essais indépendants avec probabilité de succès $p$, alors la proportion empirique $\widehat{P}$ vérifie :
$$
\mathbb{E}[\widehat{P}] = p,
\qquad
\Var(\widehat{P}) = \frac{p(1-p)}{n}.
$$
\end{Prop}

\begin{Prop}
Pour $n$ suffisamment grand, on a l'approximation normale
$$
\frac{\widehat{P}-p}{\sqrt{p(1-p)/n}} \approx \mathcal{N}(0,1).
$$
\end{Prop}

\section{Quand l'approximation normale est-elle pertinente ?}

\begin{Meth}
Dans ce cours, on utilise l'approximation normale :
\begin{itemize}
    \item pour une moyenne d'échantillon quand $n$ est suffisamment grand ;
    \item pour une proportion quand $np$ et $n(1-p)$ ne sont pas trop petits.
\end{itemize}

En pratique, dans les exercices, l'énoncé ou le contexte donne en général les conditions pour l'utiliser.
\end{Meth}

\section{Interprétation pratique}

\begin{Ex}
\textbf{Contrôle d'un lot de béton}

On mesure la résistance de $n=36$ éprouvettes. Si la moyenne vraie est $\mu$ et l'écart-type $\sigma$, alors :
$$
\overline{X} \approx \mathcal{N}\left(\mu,\frac{\sigma^2}{36}\right).
$$

L'écart-type de la moyenne vaut $\sigma/6$. La moyenne d'échantillon est donc beaucoup moins dispersée que la résistance individuelle.
\end{Ex}

\section{Erreurs fréquentes}

\begin{itemize}
    \item Confondre la dispersion d'une observation individuelle et celle d'une moyenne.
    \item Oublier le facteur $\sqrt{n}$ dans les formules.
    \item Utiliser le TCL sans justifier que l'on travaille sur un échantillon suffisamment grand.
    \item Confondre proportion empirique $\widehat{p}$ et probabilité théorique $p$.
\end{itemize}

\section{Ce qu'il faut savoir faire}

\begin{itemize}
    \item calculer une moyenne et une proportion empiriques ;
    \item interpréter un écart-type empirique ;
    \item reconnaître la loi approximative d'une moyenne d'échantillon ;
    \item reconnaître la loi approximative d'une proportion empirique ;
    \item utiliser le TCL comme outil de calcul et non comme résultat à démontrer.
\end{itemize}
